\chapter{Research, Ideation and Continuous Evaluation}{\label{ch:lit_review}}

\section{Stakeholder Research and Problem Discovery}

The project began with a user centric research phase aimed at understanding the real world
challenges within the CiPD program. Interactions were conducted with students, faculty members,
and administrative staff to identify pain points and validate potential solutions.

Students highlighted issues with manual attendance systems, describing them as inconvenient and
prone to proxy marking and lack of transparency. They also emphasized the absence of a centralized
platform for accessing schedules, resources, and tracking participation.

Faculty members pointed out difficulties in managing schedules, monitoring engagement, and
collecting structured feedback. Administrative stakeholders faced challenges in consolidating data for
reporting and decision making.

Based on these discussions, the key problem areas identified were:

\begin{itemize}
    \item Inefficient and unreliable attendance tracking
    \item Lack of a centralized integrated system
    \item Inconsistent feedback mechanisms
    \item Limited access to real time analytics
\end{itemize}

An iterative validation process was followed, where early ideas were continuously discussed with
stakeholders to ensure alignment with actual user needs.

\section{Exploration of Attendance Solutions}

Attendance tracking emerged as the most critical component, leading to the evaluation of multiple
approaches:

\begin{itemize}
    \item Bluetooth based systems relied on proximity detection but were unreliable due to device
    settings and inconsistent discoverability.
    \item RFID systems offered accuracy but required additional hardware, making them costly and
    less scalable.
    \item Face recognition approaches using computer vision were explored but faced challenges
    related to accuracy, computational overhead, and privacy concerns.
\end{itemize}

After evaluating these options, Wi Fi based detection was selected as the most practical solution. It
leverages existing infrastructure to detect connected devices and infer presence.

The choice was driven by:

\begin{itemize}
    \item No additional hardware requirement
    \item Scalability for large classrooms
    \item Passive and non intrusive operation
    \item Compatibility with existing systems
\end{itemize}

This evaluation ensured that the final approach was both feasible and scalable.

% figs/att_compare.tex — Attendance approach comparison bar chart
\begin{figure}[!ht]
\centering
\begin{tikzpicture}
\begin{axis}[
  ybar, bar width=14pt,
  symbolic x coords={Bluetooth, RFID, {Face Recog.}, {Wi-Fi}},
  xtick=data,
  ymin=0, ymax=11,
  ylabel={Suitability Score (out of 10)},
  title={Comparative Evaluation of Attendance Tracking Approaches},
  legend style={at={(0.5,-0.22)}, anchor=north, legend columns=3, font=\small},
  nodes near coords, nodes near coords align={vertical},
  every node near coord/.append style={font=\scriptsize},
  width=0.80\textwidth, height=6.5cm,
  ymajorgrids=true, grid style={dashed,gray!35},
  enlarge x limits=0.20,
]
\addplot[fill=accentblue!75, draw=accentblue]
  coordinates {(Bluetooth,4)(RFID,7)({Face Recog.},5)({Wi-Fi},9)};
\addplot[fill=successgreen!75, draw=successgreen]
  coordinates {(Bluetooth,3)(RFID,5)({Face Recog.},4)({Wi-Fi},9)};
\addplot[fill=warnorange!85, draw=warnorange]
  coordinates {(Bluetooth,6)(RFID,3)({Face Recog.},5)({Wi-Fi},9)};
\legend{Accuracy, Scalability, Cost-Effectiveness}
\end{axis}
\end{tikzpicture}
\caption{Comparative evaluation of four candidate attendance tracking approaches}
\label{fig:att_compare}
\end{figure}


\section{Network Level Experimentation and Constraints}

Following the selection of Wi Fi based attendance, experiments were conducted to test feasibility
within the institute's network. Initial trials using the existing network demonstrated that connected
devices could be detected using network scanning techniques.

However, restricted access and permission limitations prevented continuous monitoring on the
institutional network. To address this, a dedicated router was deployed within the CiPD classroom,
creating a controlled environment for reliable device detection.

Attempts were also made to enhance the router's capabilities using custom firmware (OpenWRT),
but hardware limitations restricted further customization.

These experiments highlighted key practical constraints such as network permissions, infrastructure
limitations, and hardware compatibility, which influenced the final system design.

\section{Continuous Feedback and Iterative Development}

The system was developed using an iterative approach, with regular demonstrations to stakeholders
for feedback. This allowed continuous refinement of features such as dashboards, attendance
tracking, and feedback interfaces.

Key improvements based on feedback included:

\begin{itemize}
    \item Enhanced notification mechanisms for timely updates
    \item Improved visualization of attendance data
    \item Focus on student wise analytics instead of only aggregated views
\end{itemize}

This feedback driven approach ensured that the system remained aligned with user expectations and
practical requirements while minimizing major redesign efforts.

% ─────────────────────────────────────────────────────────────────────────────
\chapter{System Overview}{\label{ch:sysoverview}}
% ─────────────────────────────────────────────────────────────────────────────

\section{Overview of CIPD 360 ERP}

CIPD 360 ERP is a web based system that provides a centralized interface for managing academic
activities. It replaces fragmented and manual processes with a cohesive platform that integrates
multiple functionalities into a single system.

The system is built with the following goals:

\begin{itemize}
    \item To provide a unified platform for all stakeholders
    \item To reduce manual intervention in routine academic processes
    \item To enable real time access to data and insights
    \item To support scalability for future expansion
\end{itemize}

By consolidating scheduling, attendance, feedback, and reporting into one system, the platform
ensures consistency and reduces redundancy across operations.

\section{Key Features Overview}

The system incorporates several key features that collectively support academic management:

\begin{itemize}
    \item \textbf{Scheduling System:}
    Enables administrators to create and manage academic sessions, with support for conflict
    detection and dynamic updates.

    \item \textbf{Automated Attendance System:}
    Uses WiFi based detection to track student presence, reducing dependency on manual
    attendance methods.

    \item \textbf{Feedback System:}
    Provides session based feedback collection with structured forms and analytical insights.

    \item \textbf{Dashboard and Reporting:}
    Offers role based dashboards with visual analytics for monitoring attendance, feedback, and
    session activity.

    \item \textbf{Content and Assignment Management:}
    Supports sharing of learning materials and management of assignments within the platform.
\end{itemize}

These features are tightly integrated, allowing smooth data flow and enabling automation across
workflows.

\section{User Roles and Responsibilities}

The system is designed around 2 primary user roles, each with distinct responsibilities and access
levels:

\begin{itemize}
    \item \textbf{Administrator:}
    Responsible for managing the overall system, including scheduling sessions, monitoring
    attendance, generating reports, and configuring system settings. As well as managing
    faculties and other teaching staff.

    \item \textbf{Student:}
    Responsible for attending sessions, accessing course materials, submitting feedback, and
    tracking personal attendance and academic progress.
\end{itemize}

\section{System Workflow}

The CIPD 360 ERP follows a structured workflow that connects different modules into a cohesive
pipeline:

\begin{enumerate}
    \item \textbf{Session Scheduling:}
    Administrators create and manage sessions, defining details such as timing, faculty, and
    venue.

    \item \textbf{Attendance Tracking:}
    During the session, the system monitors connected devices on the CiPD network and records
    presence data.

    \item \textbf{Attendance Processing:}
    After the session, detection data is processed to compute attendance records and
    participation scores.

    \item \textbf{Feedback Rollout:}
    Once a session is completed, feedback forms are automatically made available to students.

    \item \textbf{Feedback Collection and Analysis:}
    Students submit feedback, which is aggregated and analysed for insights.

    \item \textbf{Reporting and Visualization:}
    Data from all modules is consolidated and presented through dashboards and reports.
\end{enumerate}

This workflow ensures that academic activities are not isolated but part of a continuous, automated
process.

% figs/workflow.tex — System workflow cycle
\begin{figure}[!ht]
\centering
\begin{tikzpicture}[font=\small]
  \node[blknode=3.4cm] (w1) at (0,   0)    {1.\ Session\\Scheduling};
  \node[blknode=3.4cm] (w2) at (4.2, 0)    {2.\ Attendance\\Tracking};
  \node[blknode=3.4cm] (w3) at (8.4, 0)    {3.\ Attendance\\Processing};
  \node[blknode=3.4cm] (w4) at (8.4,-2.4)  {4.\ Feedback\\Rollout};
  \node[blknode=3.4cm] (w5) at (4.2,-2.4)  {5.\ Feedback\\Collection};
  \node[blknode=3.4cm] (w6) at (0,  -2.4)  {6.\ Reporting \&\\Visualization};

  \draw[myarrow] (w1.east)  -- (w2.west);
  \draw[myarrow] (w2.east)  -- (w3.west);
  \draw[myarrow] (w3.south) -- (w4.north);
  \draw[myarrow] (w4.west)  -- (w5.east);
  \draw[myarrow] (w5.west)  -- (w6.east);
  \draw[myarrow] (w6.north) -- (w1.south);
\end{tikzpicture}
\caption{End-to-end operational workflow of CIPD 360 ERP}
\label{fig:workflow}
\end{figure}


\section{System Characteristics}

The system is designed with several important characteristics:

\begin{itemize}
    \item \textbf{Modularity:}
    Each module (attendance, scheduling, feedback, etc.) operates independently but integrates
    with others.

    \item \textbf{Scalability:}
    The architecture allows the system to handle increasing numbers of users and sessions.

    \item \textbf{Automation:}
    Key workflows such as attendance computation and feedback rollout are automated.

    \item \textbf{Real Time Data Access:}
    Users have access to up to date information through dynamic dashboards.

    \item \textbf{User Centric Design:}
    Interfaces are tailored to the needs of different roles.
\end{itemize}

These characteristics ensure that the system remains robust, flexible, and efficient.
